\documentclass[12pt]{article}
\usepackage[utf8]{inputenc}

\usepackage{../mystyle}

\title{GIT 3: Quiver Representations}
\author{Joseph Sullivan}
\date{January 2026}

\DeclareMathOperator{\Kom}{Kom}
\DeclareMathOperator{\Cyl}{Cyl}
\DeclareMathOperator{\Stab}{Stab}
\DeclareMathOperator{\Slice}{Slice}
\DeclareMathOperator{\chern}{ch}
\DeclareMathOperator{\todd}{td}

\DeclareMathOperator{\Kos}{Kos}

\DeclareMathOperator{\ext}{ext}
%\DeclareMathOperator{\hom}{hom}

\newcommand{\cQ}{\mathcal{Q}}
\newcommand{\cZ}{\mathcal{Z}}
\newcommand{\LCom}{\mathcal{LC}}

\newcommand{\rddots}{\text{\reflectbox{$\ddots$}}}



\begin{document}

\maketitle

Outline
\begin{enumerate}
    \item What is a quiver representation?
    \begin{enumerate}
        \item Quiver
        \item Representations, $kQ$-modules and dimension vector
        \item A word about relations
    \end{enumerate}
    \item GIT problem
    \begin{enumerate}
        \item Space $\mathcal{R}(Q,\underline{d})$ and $\GL(\underline{d})$ group action.
        \item $G$-linearized line bundles from functionals on Grothendieck group.
        \item Sketch of Hilbert-Mumford proof/analysis
        \item Definition of $\theta$-stability a la King.
        \item Moduli functor, fact that it is a coarse moduli space up to $S$-equivalence.
    \end{enumerate}
    \item An Example
    \begin{enumerate}
        \item Generalized Kronecker quiver.
        \item Beilinson quiver.
    \end{enumerate}
\end{enumerate}

\end{document}